\chapter{Differential Forms and Stokes' Theorem}

\begin{center}\itshape
Green's theorem, the divergence theorem and the classical Stokes theorem look like three facts about three kinds of integral. Differential forms make
them one statement, $\int_M\dd\omega=\int_{\partial M}\omega$, valid in every dimension. Its first great consequence is topological: no smooth map can
retract a ball onto its boundary sphere, and therefore every continuous map of a ball into itself has a fixed point.
\end{center}

On 2 July 1850 William Thomson, the future Lord Kelvin, wrote a letter to his friend George Gabriel Stokes. In a postscript he stated, without proof,
that the flux of the curl of a vector field through a surface equals the circulation of the field around the boundary of the surface. Four years later
Stokes, who was setting the examination for Cambridge's Smith's Prize, put Thomson's statement on the paper as a question to be proved. One of the
candidates that February was a twenty-two-year-old James Clerk Maxwell, who shared the prize that year with Edward Routh. The theorem has carried Stokes'
name ever since, which is one of the better jokes in the history of mathematics: it is the most famous result that its namesake set as an exam question.

The modern form of the theorem needed a new algebra. In 1899 \'Elie Cartan introduced the exterior product of differentials, in which
$\dd x\wedge\dd y=-\dd y\wedge\dd x$, and the exterior derivative $\dd$. With them, the gradient, the curl and the divergence become a single operation, the
three classical integral theorems become one, and the proof becomes a short computation in a half-space glued together by a partition of unity.
Thirty years later Georges de Rham showed that the forms with $\dd\omega=0$, modulo those of the form $\dd\eta$, see the topology of the manifold: they count its
holes. We prove the first case of that theorem at the end of the chapter, for the plane with a point removed, where the count is the winding number.

\section{Alternating forms}

Let $V$ be a real vector space of dimension $n$. A \emph{$k$-form} on $V$ is a $k$-linear map $\alpha\colon V^k\to\R$ that is \emph{alternating}: it changes sign
when two arguments are exchanged, or equivalently vanishes when two arguments are equal. They form a vector space $\Lambda^kV^*$; $\Lambda^0V^*=\R$ and
$\Lambda^1V^*=V^*$. The model is the determinant: $n$ vectors of $\R^n$ span a parallelepiped, and its signed volume is an alternating $n$-linear function of
them. A $k$-form measures signed $k$-dimensional volumes.

Fix a basis $e_1,\dots,e_n$ of $V$ with dual basis $e^1,\dots,e^n$. For a sequence $I=(i_1,\dots,i_k)$ of indices define
\[
e^I(v_1,\dots,v_k)=\det\bigl(e^{i_m}(v_j)\bigr)_{1\le m,j\le k},
\]
the determinant of the $k\times k$ minor formed by the rows $i_1,\dots,i_k$ of the coordinates of $v_1,\dots,v_k$. It is a $k$-form; it is $0$ if $I$ has a repeated
index, and $e^{\sigma(I)}=\operatorname{sgn}(\sigma)e^I$ if the indices are permuted.

\begin{proposition}\label{prop:altbasis}
The $e^I$ with $I$ increasing form a basis of $\Lambda^kV^*$, and
\[
\alpha=\sum_{I\ \mathrm{increasing}}\alpha(e_{i_1},\dots,e_{i_k})\,e^I.
\] So $\dim\Lambda^kV^*=\binom nk$, and
$\Lambda^kV^*=0$ for $k>n$.
\end{proposition}
\begin{proof}
For increasing $I$ and $J$, $e^I(e_{j_1},\dots,e_{j_k})$ is $1$ if $I=J$ and $0$ otherwise. By multilinearity, a $k$-form is determined by its values on
sequences of basis vectors; by the alternating property, by its values on increasing ones. Both sides of the formula agree there, so they are equal, and the
same evaluation shows that the $e^I$ are independent.
\end{proof}

\begin{definition}
The \emph{exterior product} $\wedge\colon\Lambda^kV^*\times\Lambda^lV^*\to\Lambda^{k+l}V^*$ is the bilinear map with $e^I\wedge e^J=e^{IJ}$, where $IJ$ is the concatenated
sequence.
\end{definition}

This defines $\wedge$ on basis elements, and by the sign rule $e^{\sigma(I)}=\operatorname{sgn}(\sigma)e^I$ the formula $e^I\wedge e^J=e^{IJ}$ then holds for all sequences, not only
increasing ones. Two consequences are immediate. The product is associative, since $(e^I\wedge e^J)\wedge e^K=e^{IJK}=e^I\wedge(e^J\wedge e^K)$. And it is
\emph{graded commutative}: $\alpha\wedge\beta=(-1)^{kl}\beta\wedge\alpha$ for $\alpha\in\Lambda^k$, $\beta\in\Lambda^l$, because moving the $l$ indices of $J$ in front of
the $k$ indices of $I$ takes $kl$ transpositions. In particular $\lambda\wedge\lambda=0$ for a $1$-form $\lambda$.

\begin{proposition}\label{prop:wedgedet}
For $1$-forms $\lambda_1,\dots,\lambda_k$ and vectors $v_1,\dots,v_k$,
\[
(\lambda_1\wedge\dots\wedge\lambda_k)(v_1,\dots,v_k)=\det\bigl(\lambda_i(v_j)\bigr).
\]
In particular the exterior product does not depend on the basis used to define it. If $A\colon W\to V$ is linear and $A^*\alpha(w_1,\dots,w_k)=\alpha(Aw_1,\dots,Aw_k)$,
then $A^*(\alpha\wedge\beta)=A^*\alpha\wedge A^*\beta$, and for $W=V$, $A^*\omega=(\det A)\,\omega$ for every $\omega\in\Lambda^nV^*$.
\end{proposition}
\begin{proof}
Write $\lambda_i=\sum_a\lambda_{i,a}e^a$. Expanding the left side by multilinearity of $\wedge$ gives
\[\sum_{a_1,\dots,a_k}\prod_i\lambda_{i,a_i}\,e^{a_1\cdots a_k}(v_1,\dots,v_k),\]
where $e^{a_1\cdots a_k}(v)=\det(e^{a_i}(v_j))_{i,j}$. Expanding the right side by
multilinearity of the determinant in its rows, $\det(\sum_a\lambda_{i,a}e^a(v_j))_{i,j}$, gives the same sum. The formula involves no basis, and since every
$k$-form is a sum of products of $1$-forms (the $e^I$ are such products), $\wedge$ is basis-free. For the pullback, $A^*(\lambda_1\wedge\dots\wedge\lambda_k)=A^*\lambda_1\wedge\dots\wedge A^*\lambda_k$
by the determinant formula, and the general case follows by bilinearity. Finally $A^*(e^1\wedge\dots\wedge e^n)(e_1,\dots,e_n)=\det(e^i(Ae_j))=\det A$, and $\Lambda^nV^*$ is
one-dimensional.
\end{proof}

\section{Differential forms and the exterior derivative}

\begin{definition}
A \emph{differential $k$-form} on a manifold $M$ is a choice of $\omega_p\in\Lambda^kT_p^*M$ for every $p$, smooth in the sense that in every chart
\[
\omega=\sum_{I\ \mathrm{increasing}}a_I\,\dd x^I,\qquad\dd x^I=\dd x^{i_1}\wedge\dots\wedge\dd x^{i_k},
\]
with smooth coefficients $a_I$, where $\dd x^1,\dots,\dd x^n$ is the basis of $T_p^*M$ dual to the $\frac\partial{\partial x^i}$. The space of $k$-forms is $\Omega^k(M)$; a
$0$-form is a smooth function. For a smooth $F\colon M\to N$ and $\omega\in\Omega^k(N)$, the \emph{pullback} $F^*\omega\in\Omega^k(M)$ is
$(F^*\omega)_p(v_1,\dots,v_k)=\omega_{F(p)}(dF_pv_1,\dots,dF_pv_k)$.
\end{definition}

The differential of a function $f$ is the $1$-form $\dd f_p=df_p$, with $\dd f=\sum_i\partial_if\,\dd x^i$ in a chart, which explains the notation $\dd x^i$ for the
differential of the coordinate function $x^i$. By Proposition~\ref{prop:wedgedet}, $F^*(\alpha\wedge\beta)=F^*\alpha\wedge F^*\beta$, $F^*f=f\circ F$ for functions, and the chain
rule says $F^*\dd f=\dd(f\circ F)$. So in coordinates pullback is substitution: $F^*(a\,\dd y^{j_1}\wedge\dots)=(a\circ F)\,\dd F^{j_1}\wedge\dots$, where
$F^j=y^j\circ F$.

\begin{theorem}[exterior derivative]\label{thm:extd}
There is a unique family of linear maps $\dd\colon\Omega^k(M)\to\Omega^{k+1}(M)$ such that $\dd f$ is the differential for functions, $\dd\circ\dd=0$, and
\[
\dd(\alpha\wedge\beta)=\dd\alpha\wedge\beta+(-1)^k\alpha\wedge\dd\beta\qquad(\alpha\in\Omega^k).
\]
In a chart, $\dd\bigl(\sum_Ia_I\dd x^I\bigr)=\sum_I\dd a_I\wedge\dd x^I$. Moreover $F^*\dd\omega=\dd F^*\omega$ for every smooth $F$.
\end{theorem}
\begin{proof}
\emph{On an open set $U\subset\R^n$.} Define $\dd$ by the formula. It is linear and agrees with the differential on functions. Leibniz: by linearity it
suffices to take $\alpha=a\,\dd x^I$, $\beta=b\,\dd x^J$; then $\alpha\wedge\beta=ab\,\dd x^I\wedge\dd x^J$ and
\[
\dd(\alpha\wedge\beta)=(b\,\dd a+a\,\dd b)\wedge\dd x^I\wedge\dd x^J=\dd\alpha\wedge\beta+(-1)^ka\,\dd x^I\wedge\dd b\wedge\dd x^J,
\]
moving the $1$-form $\dd b$ past the $k$-form $\dd x^I$. The formula also holds when $I$ is not increasing, by the sign rule. For $\dd\dd=0$: on functions,
$\dd\dd f=\sum_{i,j}\partial_j\partial_if\,\dd x^j\wedge\dd x^i=0$, since the coefficients are symmetric in $(i,j)$ and $\dd x^j\wedge\dd x^i$ is antisymmetric; then
$\dd\dd(a\,\dd x^I)=\dd(\dd a\wedge\dd x^I)=\dd\dd a\wedge\dd x^I-\dd a\wedge\dd(\dd x^I)=0$, as $\dd(\dd x^I)=\dd(1)\wedge\dd x^I=0$. Uniqueness: any operator with the three properties
satisfies $\dd(\dd x^I)=0$ (by Leibniz and $\dd\dd x^i=0$), hence $\dd(a\,\dd x^I)=\dd a\wedge\dd x^I$ by Leibniz.

\emph{Pullback.} For $F\colon U\to V$ between open sets of Euclidean spaces and $\omega=a\,\dd y^{j_1}\wedge\dots\wedge\dd y^{j_k}$,
$F^*\omega=(a\circ F)\,\dd F^{j_1}\wedge\dots\wedge\dd F^{j_k}$, so by Leibniz and $\dd\dd=0$,
\[
\dd F^*\omega=\dd(a\circ F)\wedge\dd F^{j_1}\wedge\dots=F^*\dd a\wedge F^*\dd y^{j_1}\wedge\dots=F^*(\dd\omega).
\]

\emph{On a manifold.} Define $\dd\omega$ in each chart by the formula. On the overlap of two charts the results agree, because the transition map commutes
with $\dd$. The properties are local, so they hold on $M$, and so does $F^*\dd=\dd F^*$. Uniqueness on $M$: the three properties force $\dd$ to be local (if
$\omega=0$ near $p$, take a bump $\chi$ equal to $1$ near $p$ with $\chi\omega=0$; then $0=\dd(\chi\omega)=\dd\chi\wedge\omega+\chi\,\dd\omega$, which is $\dd\omega$ at $p$), and locally
it is given by the formula.
\end{proof}

\begin{example}[grad, curl, div]\label{ex:gradcurldiv}
In $\R^3$ a function $f$, a vector field $F=(P,Q,R)$ and a function $\rho$ correspond to the forms $f$, $\omega_F=P\dd x+Q\dd y+R\dd z$,
$\eta_F=P\dd y\wedge\dd z+Q\dd z\wedge\dd x+R\dd x\wedge\dd y$ and $\rho\,\dd x\wedge\dd y\wedge\dd z$. Then $\dd f=\omega_{\nabla f}$, $\dd\omega_F=\eta_{\nabla\times F}$ and
$\dd\eta_F=(\nabla\cdot F)\dd x\wedge\dd y\wedge\dd z$, and $\dd\dd=0$ becomes $\nabla\times\nabla f=0$ and $\nabla\cdot(\nabla\times F)=0$ (Exercise~\ref{ex:gradcurl}).
\end{example}

\section{Orientation and integration}

\begin{definition}
Two charts are \emph{orientation-compatible} if their transition maps have positive Jacobian determinant. An \emph{orientation} of $M$ is a maximal atlas of
pairwise orientation-compatible charts, and $M$ is \emph{orientable} if it has one.
\end{definition}

\begin{proposition}\label{prop:volumeform}
A manifold of dimension $n\ge1$ is orientable if and only if it carries a nowhere-vanishing $n$-form. Such a form $\mu$ determines the orientation of the
charts in which $\mu=a\,\dd x^1\wedge\dots\wedge\dd x^n$ with $a>0$.
\end{proposition}
\begin{proof}
If $\mu$ vanishes nowhere, then in any chart on a connected domain $a$ has constant sign, and composing with a reflection of one coordinate we may make
it positive. For two such charts, $\mu=a\,\dd x^1\wedge\dots=b\,\dd y^1\wedge\dots$ and Proposition~\ref{prop:wedgedet} gives $a=b\det\bigl(\frac{\partial y^i}{\partial x^j}\bigr)$, so the
transition has positive Jacobian. Conversely, given an oriented atlas, take a partition of unity $(\rho_j)$ with each $\operatorname{supp}\rho_j$ in an oriented chart
domain $U_j$ (Theorem~\ref{thm:partition}), and $\mu=\sum_j\rho_j\,\dd x_j^1\wedge\dots\wedge\dd x_j^n$. In any oriented chart, each term is a non-negative multiple of
$\dd x^1\wedge\dots\wedge\dd x^n$, and at each point at least one is positive.
\end{proof}

The sphere is orientable, by the form of Example~\ref{ex:sigma} below; the M\"obius band and $\R P^2$ are not (Exercise~\ref{ex:mobius}).

\begin{definition}
Let $M$ be oriented and $\omega$ an $n$-form with compact support. If the support lies in the domain of an oriented chart $\varphi$, write
$(\varphi^{-1})^*\omega=a\,\dd x^1\wedge\dots\wedge\dd x^n$ and put $\int_M\omega=\int_{\varphi(U)}a\dd\lambda$. In general choose a partition of unity $(\rho_j)$ subordinate to a cover by
oriented charts and put $\int_M\omega=\sum_j\int_M\rho_j\omega$, a finite sum.
\end{definition}

\begin{proposition}\label{prop:intwelldef}
The integral is well defined and linear. If $F\colon N\to M$ is a diffeomorphism that preserves orientation, $\int_NF^*\omega=\int_M\omega$; if it reverses it, the sign
changes.
\end{proposition}
\begin{proof}
For a form supported in the domains of two oriented charts $\varphi$ and $\psi$, let $T=\psi\circ\varphi^{-1}$ and $(\psi^{-1})^*\omega=b\,\dd y^1\wedge\dots\wedge\dd y^n$. Then
$(\varphi^{-1})^*\omega=T^*\bigl(b\,\dd y^1\wedge\dots\bigr)=(b\circ T)\det(DT)\,\dd x^1\wedge\dots$ by Proposition~\ref{prop:wedgedet}, and since $\det DT>0$ the change of variables
formula (Volume~6, Theorem~8.8) gives $\int(b\circ T)\det DT\dd\lambda=\int b\dd\lambda$. The sum over a partition of unity is finite, because the compact support
meets only finitely many supports. If $(\sigma_i)$ is another partition,
$\sum_j\int\rho_j\omega=\sum_j\sum_i\int\sigma_i\rho_j\omega=\sum_i\int\sigma_i\omega$ by linearity of the one-chart integral. The statement about $F$ is the same computation
with $\det DT$ replaced by the Jacobian of $F$ in charts, positive or negative.
\end{proof}

A $k$-form can be integrated over any oriented $k$-dimensional submanifold, compact or with compact support in it, by pulling it back. For $k=1$ and a curve
$\gamma\colon[a,b]\to M$, $\int_\gamma\omega=\int_a^b\gamma^*\omega=\int_a^b\omega_{\gamma(t)}(\gamma'(t))\dd t$, and for $\omega=\dd f$ this is $f(\gamma(b))-f(\gamma(a))$, the
fundamental theorem of calculus.

\section{Stokes' theorem}

\begin{definition}
Let $H^n=\{x\in\R^n:x_1\le0\}$. A \emph{manifold with boundary} is defined like a manifold, with charts taking values in open subsets of $H^n$; a map on a
subset of $H^n$ is smooth if it extends to a smooth map on an open set of $\R^n$. The \emph{boundary} $\partial M$ is the set of points sent to
$\partial H^n=\{x_1=0\}$ by a chart.
\end{definition}

The boundary is well defined: if a transition map $T$ sent a point $p$ with $p_1<0$ to a point $q$ with $q_1=0$, its derivative at $p$ would be invertible, and by
the inverse function theorem $T$ would map a neighbourhood of $p$ onto an open set of $\R^n$ containing $q$, which is not contained in $H^n$. The
restrictions of the charts to $\{x_1=0\}$ make $\partial M$ a manifold of dimension $n-1$ without boundary. A closed ball $\bar B\subset\R^n$ is a manifold with
boundary $S^{n-1}$; so is a compact region of the plane bounded by finitely many smooth curves (Figure~\ref{fig:boundary}).

\begin{proposition}[boundary orientation]\label{prop:bdryorient}
Let $M$ be oriented and $n\ge2$. The charts $(x^2,\dots,x^n)$ of $\partial M$ obtained by restricting oriented charts of $M$ are orientation-compatible, and define the
\emph{induced orientation} on $\partial M$: ``outward normal first''.
\end{proposition}
\begin{proof}
A transition map $T$ between two oriented charts maps $\{x_1=0\}$ to itself and $\{x_1<0\}$ to itself. So $T_1(0,x')=0$, whence $\partial_jT_1=0$ for $j\ge2$ at
boundary points, and $\partial_1T_1\ge0$ there. The Jacobian matrix is block triangular, $\det DT=\partial_1T_1\cdot\det D'T'$, where $T'$ is the transition between the boundary
charts. Since $\det DT>0$, both factors are positive.
\end{proof}

In $H^n$ the outward direction is $e_1$, and $(e_1,e_2,\dots,e_n)$ is positive: the boundary is oriented so that the outward normal followed by a positive basis
of the boundary is a positive basis of $M$. For a plane region this says that the outer boundary turns anticlockwise and the boundaries of holes clockwise,
with the region on the left (Figure~\ref{fig:boundary}). For $n=1$ we orient a boundary point $+$ or $-$ according as an oriented chart maps it to the
boundary of $H^1=\left]-\infty,0\right]$ or of $[0,\infty[$: the end of an interval is $+$, the beginning $-$.

\begin{figure}[ht]
\centering
\fig{ch03-boundary}
\caption{A compact plane region $M$ with one hole, a manifold with boundary. With the standard orientation of the plane, the induced orientation of
$\partial M$ runs anticlockwise on the outer curve and clockwise on the inner one: in both cases the outward normal (orange) followed by the direction of the
boundary is a positive basis.}
\label{fig:boundary}
\end{figure}

\begin{theorem}[Stokes]\label{thm:stokes}
Let $M$ be an oriented $n$-manifold with boundary, with $\partial M$ carrying the induced orientation, and $\omega$ an $(n-1)$-form with compact support. Then
\[
\int_M\dd\omega=\int_{\partial M}\omega.
\]
\end{theorem}
\begin{proof}
Both sides are linear in $\omega$, and a partition of unity writes $\omega$ as a finite sum of forms supported in single oriented charts. So it suffices to
treat $\omega$ supported in the domain of one oriented chart, and by Proposition~\ref{prop:intwelldef} and $F^*\dd=\dd F^*$ we may compute in the chart: $\omega$ is an
$(n-1)$-form with compact support in an open set of $H^n$, which we extend by $0$ to $H^n$. Write
\[
\omega=\sum_{j=1}^na_j\,\dd x^1\wedge\dots\wedge\widehat{\dd x^j}\wedge\dots\wedge\dd x^n,\qquad\dd\omega=\sum_j(-1)^{j-1}\partial_ja_j\,\dd x^1\wedge\dots\wedge\dd x^n,
\]
where the hat means omission and the sign comes from moving $\dd x^j$ to its place. Integrate each term by Fubini (Volume~6). For $j\ge2$, the integral of
$\partial_ja_j$ over the line in the variable $x_j$ is $0$, since $a_j$ has compact support. For $j=1$,
\[
\int_{H^n}\partial_1a_1\dd\lambda=\int_{\R^{n-1}}\Bigl(\int_{-\infty}^0\partial_1a_1(x_1,x')\dd x_1\Bigr)\dd x'=\int_{\R^{n-1}}a_1(0,x')\dd x'.
\]
On the boundary $\{x_1=0\}$ the form $\dd x^1$ restricts to $0$, so only the term $j=1$ of $\omega$ survives, and in the positive chart $(x^2,\dots,x^n)$ of $\partial H^n$
its integral is exactly $\int a_1(0,x')\dd x'$. If the chart does not meet the boundary ($\omega$ supported in $\{x_1<0\}$), the same computation gives $0$ on
both sides. For $n=1$, the theorem is the fundamental theorem of calculus with the sign convention above.
\end{proof}

\begin{beauty}
All the difficulty of the classical theorems has gone into the definitions. The proof itself is the fundamental theorem of calculus in one variable,
applied once, in a half-space; everything else is bookkeeping by partitions of unity, and the cancellations that make the boundary of a boundary empty are
built into the algebra $\dd x\wedge\dd y=-\dd y\wedge\dd x$.
\end{beauty}

\begin{example}[the volume form of the sphere]\label{ex:sigma}
On $\R^n$ let $\sigma=\sum_{i=1}^n(-1)^{i-1}x_i\,\dd x^1\wedge\dots\wedge\widehat{\dd x^i}\wedge\dots\wedge\dd x^n$. Then $\dd\sigma=n\,\dd x^1\wedge\dots\wedge\dd x^n$, so by Stokes on the closed unit
ball,
\[
\int_{S^{n-1}}\sigma=n\,\lambda(B^n)>0.
\]
At a point $x\in S^{n-1}$, $\sigma_x(v_2,\dots,v_n)=\det(x,v_2,\dots,v_n)$ (expand along the first column), which is the volume spanned by the unit outward normal $x$ and
the $v_i$: the restriction of $\sigma$ is the area form of the sphere, it vanishes nowhere, and $S^{n-1}$ is orientable. For $n=2$, $\sigma=x\dd y-y\dd x$ and
$\int_{S^1}\sigma=2\pi$; for $n=3$, the area of $S^2$ is $3\cdot\frac43\pi=4\pi$.
\end{example}

\subsection*{Brouwer's fixed point theorem in every dimension}

Volume~7 proved Brouwer's theorem for the disc with winding numbers. Stokes' theorem gives it in every dimension in a few lines.

\begin{theorem}[no smooth retraction]\label{thm:noretraction}
There is no smooth map $r\colon\bar B^n\to S^{n-1}$ with $r(x)=x$ for $x\in S^{n-1}$.
\end{theorem}
\begin{proof}
Suppose there is one, and let $\sigma$ be the form of Example~\ref{ex:sigma}, restricted to $S^{n-1}$. Since $r$ is the identity on the boundary,
\[
0<\int_{S^{n-1}}\sigma=\int_{S^{n-1}}r^*\sigma=\int_{\bar B^n}\dd(r^*\sigma)=\int_{\bar B^n}r^*(\dd\sigma)=0,
\]
because $\dd\sigma$, computed on the $(n-1)$-dimensional manifold $S^{n-1}$, is an $n$-form on it, hence zero (Proposition~\ref{prop:altbasis}).
\end{proof}

\begin{theorem}[Brouwer]\label{thm:brouwern}
Every continuous map $f\colon\bar B^n\to\bar B^n$ has a fixed point.
\end{theorem}
\begin{proof}
Suppose $f(x)\ne x$ for all $x$; then $m=\min|f(x)-x|>0$. By the Stone--Weierstrass theorem (Volume~6, Theorem~4.19) choose a polynomial map $P$ with
$|P-f|<\frac m4$ on $\bar B^n$, and let $g=\frac P{1+m/4}$. Then $|g|<1$ on $\bar B^n$, and $|g-f|\le|P-f|+\frac{m/4}{1+m/4}|P|<\frac m4+\frac m4$, so $g(x)\ne x$. The map $g$
is smooth and has no fixed point; follow the ray from $g(x)$ through $x$ until it meets the sphere:
\[
r(x)=x+t(x)\,u(x),\quad u=x-g(x),\quad t=\frac{-\langle x,u\rangle+\sqrt{\langle x,u\rangle^2+|u|^2(1-|x|^2)}}{|u|^2},
\]
the non-negative root of $|x+tu|^2=1$. The square root is of a positive quantity on a neighbourhood of $\bar B^n$: for $|x|<1$ the second term is positive, and for
$|x|=1$, $\langle x,u\rangle=1-\langle x,g(x)\rangle>0$ since $|g(x)|<1$. So $r$ is smooth, maps $\bar B^n$ to $S^{n-1}$, and for $|x|=1$ we get $t=0$, $r(x)=x$. This contradicts
Theorem~\ref{thm:noretraction}.
\end{proof}

\section{De Rham cohomology}

A form is \emph{closed} if $\dd\omega=0$ and \emph{exact} if $\omega=\dd\eta$. Exact forms are closed, since $\dd\dd=0$. The question of the converse is the question of
whether a vector field with zero curl is a gradient, and its answer depends on the shape of the domain.

\begin{definition}
The \emph{de Rham cohomology} of $M$ in degree $k$ is the quotient vector space
\[H^k(M)=\frac{\{\text{closed }k\text{-forms}\}}{\{\text{exact }k\text{-forms}\}}.\] A smooth $F\colon M\to N$ induces
$F^*\colon H^k(N)\to H^k(M)$, since $F^*$ commutes with $\dd$.
\end{definition}

$H^0(M)$ is the space of locally constant functions, of dimension the number of connected components. The higher groups are computed by the following
formula, which says that a smooth deformation changes forms only by exact ones.

\begin{theorem}[homotopy formula]\label{thm:homotopy}
Let $h\colon\R\times M\to N$ be smooth and $h_t=h(t,\cdot)$. There are linear maps $K\colon\Omega^k(N)\to\Omega^{k-1}(M)$ with
\[
h_1^*\omega-h_0^*\omega=\dd K\omega+K\dd\omega.
\]
In particular smoothly homotopic maps induce the same map on cohomology.
\end{theorem}
\begin{proof}
Every form on $\R\times M$ can be written uniquely as $\alpha_t+\dd t\wedge\beta_t$, where $\alpha_t$ and $\beta_t$ are forms on $M$ depending smoothly on $t$ (in a chart
$(t,x)$, collect the terms containing $\dd t$). Write $h^*\omega=\alpha_t+\dd t\wedge\beta_t$ and define $K\omega=\int_0^1\beta_t\dd t$, integrating the coefficients. Restricting to
the slice $\{t=c\}$ kills $\dd t$, so $\alpha_c=h_c^*\omega$. In a chart, $\dd(\alpha_t)=\dd t\wedge\partial_t\alpha_t+\dd_M\alpha_t$ and $\dd(\dd t\wedge\beta_t)=-\dd t\wedge\dd_M\beta_t$, where $\dd_M$ differentiates
in the $M$ variables only. Hence
\[
h^*\dd\omega=\dd h^*\omega=\dd_M\alpha_t+\dd t\wedge(\partial_t\alpha_t-\dd_M\beta_t),
\]
and by definition of $K$,
\[
K\dd\omega=\int_0^1(\partial_t\alpha_t-\dd_M\beta_t)\dd t=\alpha_1-\alpha_0-\dd K\omega,
\]
where we differentiated under the integral sign. This is the formula. If $\omega$ is closed, $h_1^*\omega-h_0^*\omega=\dd K\omega$ is exact.
\end{proof}

\begin{corollary}[Poincar\'e lemma]\label{cor:poincare}
On a star-shaped open set $U\subset\R^n$, every closed $k$-form with $k\ge1$ is exact: $H^k(U)=0$.
\end{corollary}
\begin{proof}
Say $U$ is star-shaped about $0$, and let $h(t,x)=\beta(t)x$ with $\beta$ the smooth function of Lemma~\ref{lem:bump}, which takes values in $[0,1]$ and equals $0$ at
$t=0$ and $1$ at $t=1$; it maps $\R\times U$ to $U$. Then $h_1=\mathrm{id}$ and $h_0$ is constant, so $h_0^*\omega=0$ for $k\ge1$ and $\omega=\dd K\omega$ for closed $\omega$.
\end{proof}

\subsection*{The punctured plane and the winding number}

On $\R^2\setminus\{0\}$, the $1$-form
\[
\dd\theta=\frac{x\dd y-y\dd x}{x^2+y^2}
\]
is closed (a direct computation), and its integral over the unit circle is $2\pi$ (it restricts to $\sigma$ there, Example~\ref{ex:sigma}). So it is not exact,
for the integral of $\dd f$ over a closed curve is $0$. The notation is suggestive: locally $\dd\theta$ is the differential of an angle function, but there is no
global one.

\begin{theorem}[de Rham's theorem for the punctured plane]\label{thm:h1punctured}
A closed $1$-form $\omega$ on $\R^2\setminus\{0\}$ is exact if and only if $\int_{S^1}\omega=0$. Hence $H^1(\R^2\setminus\{0\})$ is one-dimensional, spanned by the class of
$\dd\theta$, and the isomorphism with $\R$ is $\omega\mapsto\frac1{2\pi}\int_{S^1}\omega$.
\end{theorem}
\begin{proof}
Let $P(r,\phi)=(r\cos\phi,r\sin\phi)$ on the half-plane $\Pi=\left]0,\infty\right[\times\R$, which is convex. The form $P^*\omega=a\,\dd r+b\,\dd\phi$ is closed, with coefficients
$2\pi$-periodic in $\phi$ since $P$ is. By Corollary~\ref{cor:poincare} there is $g$ on $\Pi$ with $\dd g=P^*\omega$. Let $c(r,\phi)=g(r,\phi+2\pi)-g(r,\phi)$. Its partial derivatives
are $a(r,\phi+2\pi)-a(r,\phi)=0$ and $b(r,\phi+2\pi)-b(r,\phi)=0$, so $c$ is a constant, and taking $r=1$ it equals $\int_0^{2\pi}b(1,\phi)\dd\phi=\int_{S^1}\omega$. If this is $0$, $g$ is
$2\pi$-periodic in $\phi$, so $g=f\circ P$ for a function $f$ on $\R^2\setminus\{0\}$, smooth because $P$ is a local diffeomorphism, and $P^*\dd f=\dd g=P^*\omega$ gives $\dd f=\omega$. For
the second statement, $\omega-\frac c{2\pi}\dd\theta$ has integral $0$ over $S^1$, hence is exact.
\end{proof}

\begin{corollary}[the winding number as an integral]\label{cor:windingintegral}
For a closed piecewise $C^1$ curve $\gamma$ in $\R^2\setminus\{a\}$, the winding number of Volume~7 (Definition after Lemma~9.1) is
$n(\gamma,a)=\frac1{2\pi}\int_\gamma\dd\theta_a$, where $\dd\theta_a$ is $\dd\theta$ translated to $a$.
\end{corollary}
\begin{proof}
Take $a=0$. With $\theta(t)$ a continuous argument of $\gamma(t)$ (Volume~7, Lemma~9.1), on each short piece of $\gamma$ the function $\theta$ coincides with a smooth local angle
function $\vartheta$ composed with $\gamma$, and $\dd\vartheta=\dd\theta$. So $\int_\gamma\dd\theta$ is the sum of the increments of $\theta$ over the pieces, which is $\theta(\beta)-\theta(\alpha)=2\pi\,n(\gamma,0)$.
\end{proof}

\begin{remark}
Theorem~\ref{thm:h1punctured} is de Rham's theorem for the punctured plane, and only for that case: the one hole is detected by the one closed form that winds around it. For a general compact manifold, de Rham's theorem identifies $H^k(M)$ with the dual of the singular homology $H_k(M;\R)$ of Chapter~6. That identification is not proved in this book.
\end{remark}

\section{Exercises}

\begin{exercise}[$\star$]\label{ex:gradcurl}
\emph{Vector calculus.} Verify the three formulas of Example~\ref{ex:gradcurldiv} and deduce $\nabla\times(\nabla f)=0$ and $\nabla\cdot(\nabla\times F)=0$. Then deduce from
Theorem~\ref{thm:stokes}: (a) Green's theorem $\int_M(\partial_xQ-\partial_yP)\dd x\dd y=\int_{\partial M}P\dd x+Q\dd y$ for a compact plane region; (b) the divergence theorem
$\int_M\nabla\cdot F\dd V=\int_{\partial M}\eta_F$ for a compact region of $\R^3$; (c) the classical Stokes theorem for a compact oriented surface $\Sigma\subset\R^3$ with
boundary: $\int_\Sigma\eta_{\nabla\times F}=\int_{\partial\Sigma}\omega_F$.
\end{exercise}

\begin{exercise}[$\star$]\label{ex:shoelace}
\emph{The shoelace formula.} (a) Show that the area of a compact plane region $M$ is $\frac12\int_{\partial M}(x\dd y-y\dd x)$. (b) Deduce that a simple polygon with
vertices $(x_1,y_1),\dots,(x_m,y_m)$ listed anticlockwise has area $\frac12\sum_i(x_iy_{i+1}-x_{i+1}y_i)$ (indices modulo $m$), and interpret each term as the signed area of
a triangle with a vertex at the origin (Figure~\ref{fig:shoelace}). (Polygons have corners and are not manifolds with boundary; smooth the corners, or apply
(a) to each triangle.)
\end{exercise}

\begin{exercise}[$\star$]\label{ex:windingcompute}
\emph{Winding numbers by integration.} (a) Check that $\dd\theta$ is closed. (b) Compute $\frac1{2\pi}\int_\gamma\dd\theta$ for $\gamma(t)=(\cos t,\sin t)$ on $[0,2k\pi]$ and for
$\gamma(t)=(\cos2t+\frac12\cos t,\ \sin2t+\frac12\sin t)$ on $[0,2\pi]$. (c) Show that $\frac1{2\pi}\int_\gamma\dd\theta$ does not change if $\gamma$ is deformed
by a smooth homotopy of closed curves in $\R^2\setminus\{0\}$.
\end{exercise}

\begin{exercise}[$\star$]\label{ex:toruscohomology}
\emph{Holes of the torus.} On $T^2=\R^2/\Z^2$ the forms $\dd\theta_1$ and $\dd\theta_2$ are well defined. (a) Show they are closed and not exact. (b) Show that their classes
are independent in $H^1(T^2)$ by integrating over the two circles $\theta_2=0$ and $\theta_1=0$. (c) Show that $\dd\theta_1\wedge\dd\theta_2$ is closed and not exact. (In fact
$H^1(T^2)=\R^2$ and $H^2(T^2)=\R$.)
\end{exercise}

\begin{exercise}[$\star\star$]\label{ex:mobius}
\emph{Non-orientable surfaces.} (a) Let $A$ be the antipodal map of $S^2$. Show that $A^*\sigma=-\sigma$ for the area form $\sigma$ of Example~\ref{ex:sigma} with $n=3$. (b) Deduce that
$\R P^2$ is not orientable: a nowhere-vanishing $2$-form $\mu$ on $\R P^2$ would pull back to $S^2$ as $f\sigma$ with $f$ never zero and $f\circ A=-f$. (c) Why does the same argument
show that $\R P^n$ is orientable when $n$ is odd?
\end{exercise}

\begin{exercise}[$\star\star$]\label{ex:spherevolume}
\emph{Volumes of balls and spheres.} Let $\omega_n=\lambda(B^n)$ and $s_{n-1}=\int_{S^{n-1}}\sigma$ the area of the unit sphere. (a) Show $s_{n-1}=n\omega_n$. (b) Using
$\int_{\R^n}e^{-|x|^2}\dd x=\pi^{n/2}$ (Volume~6) and polar coordinates, show $s_{n-1}=\frac{2\pi^{n/2}}{\Gamma(n/2)}$. (c) Deduce $\omega_n=\frac{\pi^{n/2}}{\Gamma(\frac n2+1)}$, and that
$\omega_n\to0$: in high dimensions the unit ball has almost no volume.
\end{exercise}

\section{Hints}

\begin{hint}{ex:gradcurl}
For (b) and (c), the forms $\eta_F$ and $\omega_F$ and Theorem~\ref{thm:stokes} with $n=3$ and $n=2$; for (a), $\omega=P\dd x+Q\dd y$.
\end{hint}
\begin{hint}{ex:shoelace}
$\dd(x\dd y-y\dd x)=2\dd x\wedge\dd y$. On the segment from $p$ to $q$, $\int(x\dd y-y\dd x)=\det(p,q)$.
\end{hint}
\begin{hint}{ex:windingcompute}
For (b), count how often the curve turns around $0$ by following the argument; for (c), Theorem~\ref{thm:homotopy} with $M$ a circle.
\end{hint}
\begin{hint}{ex:toruscohomology}
The integral of an exact form over a closed curve is $0$, and the integral of an exact $2$-form over the compact surface $T^2$ is $0$ by Stokes.
\end{hint}
\begin{hint}{ex:mobius}
$A$ is the restriction of $-I$, and $(-I)^*$ multiplies $n$-forms by $(-1)^n$; $\sigma$ has degree $n-1$ but coefficients of degree $1$.
\end{hint}
\begin{hint}{ex:spherevolume}
In polar coordinates $\dd x=r^{n-1}\dd r\,\dd S$, and $\int_0^\infty e^{-r^2}r^{n-1}\dd r=\frac12\Gamma(\frac n2)$.
\end{hint}

\section{Solutions}

\begin{solution}{ex:gradcurl}
\idea Translate each vector field into the form that has the same integrals, and let $\dd$ do the differentiation.
\writeup $\dd f=f_x\dd x+f_y\dd y+f_z\dd z=\omega_{\nabla f}$. For $\omega_F$, $\dd(P\dd x)=P_y\dd y\wedge\dd x+P_z\dd z\wedge\dd x$, and adding the similar terms for $Q$ and $R$ and collecting,
$\dd\omega_F=(R_y-Q_z)\dd y\wedge\dd z+(P_z-R_x)\dd z\wedge\dd x+(Q_x-P_y)\dd x\wedge\dd y=\eta_{\nabla\times F}$. For $\eta_F$, only the derivative in the missing variable survives:
$\dd\eta_F=(P_x+Q_y+R_z)\dd x\wedge\dd y\wedge\dd z$. Then $\nabla\times\nabla f=0$ is $\dd\dd f=0$, and $\nabla\cdot(\nabla\times F)=0$ is $\dd\dd\omega_F=0$. (a) $\dd(P\dd x+Q\dd y)=(Q_x-P_y)\dd x\wedge\dd y$,
and Stokes applies. (b) Stokes with $\omega=\eta_F$ on a compact region $M\subset\R^3$; $\eta_F$ restricted to $\partial M$ is $\langle F,\nu\rangle\dd S$, the flux. (c) Stokes with
$\omega=\omega_F$ on the surface $\Sigma$, whose boundary curve carries the induced orientation; $\int_\Sigma\eta_{\nabla\times F}$ is the flux of the curl. This is
Thomson's postscript and Stokes' exam question, in one line.
\end{solution}

\begin{solution}{ex:shoelace}
\idea Area is the integral of $\dd x\wedge\dd y$, which is $\dd$ of several $1$-forms; the symmetric choice $\frac12(x\dd y-y\dd x)$ gives a formula with a pleasant geometric meaning.
\writeup (a) $\dd\bigl(\frac12(x\dd y-y\dd x)\bigr)=\frac12(\dd x\wedge\dd y-\dd y\wedge\dd x)=\dd x\wedge\dd y$, and Theorem~\ref{thm:stokes} gives
$\lambda(M)=\int_M\dd x\wedge\dd y=\frac12\int_{\partial M}(x\dd y-y\dd x)$. (b) On the segment $\gamma(t)=p+t(q-p)$, $x\dd y-y\dd x$ pulls back to
$(p_1+t(q_1-p_1))(q_2-p_2)-(p_2+t(q_2-p_2))(q_1-p_1)=p_1q_2-p_2q_1$, constant, so its integral is $\det(p,q)$. Summing over the edges gives the formula. The
number $\frac12\det(p,q)$ is the signed area of the triangle $0pq$, positive if the edge turns anticlockwise around the origin. The area of the polygon is the sum
of these signed areas: triangles outside the polygon are counted once positively and once negatively and cancel (Figure~\ref{fig:shoelace}). To justify
(a) for polygons, apply it to a triangle (for which the formula can also be checked directly) and decompose the polygon into triangles: the interior
edges are traversed twice in opposite directions and cancel.
\end{solution}

\begin{figure}[ht]
\centering
\fig{ch03-shoelace}
\caption{The shoelace formula. Each edge $pq$ of the polygon, traversed anticlockwise, contributes the signed area $\frac12\det(p,q)$ of the triangle $0pq$:
positive (green) when the edge turns anticlockwise around the origin, negative (red) otherwise. Where triangles of both signs overlap they cancel, and what
remains is the polygon.}
\label{fig:shoelace}
\end{figure}

\begin{solution}{ex:windingcompute}
\idea The integral of $\dd\theta$ is the total change of angle, so computing it is counting turns, and deforming the curve cannot change an integer that varies continuously.
\writeup (a) With $\rho=x^2+y^2$, $\dd\bigl(\frac x\rho\bigr)\wedge\dd y-\dd\bigl(\frac y\rho\bigr)\wedge\dd x=\bigl(\partial_x\frac x\rho+\partial_y\frac y\rho\bigr)\dd x\wedge\dd y=\frac{(\rho-2x^2)+(\rho-2y^2)}{\rho^2}\dd x\wedge\dd y=0$.
(b) On the circle $\gamma^*\dd\theta=\dd t$, so the integral over $[0,2k\pi]$ is $2k\pi$ and the winding number is $k$. For the second curve, in complex notation $\gamma(t)=e^{2it}\bigl(1+\frac12e^{-it}\bigr)$. The factor
$1+\frac12e^{-it}$ stays in the disc $D(1,\frac12)$, inside the right half-plane, where $\operatorname{Im}\operatorname{Log}$ is a continuous argument; its argument therefore
returns to its initial value. So a continuous argument of $\gamma$ increases by exactly $2\cdot2\pi$, the integral is $4\pi$ by Corollary~\ref{cor:windingintegral}, and the
winding number is $2$: the small loop of the curve also goes around the origin. (This is Volume~7's ``dog on a leash'': the term $\frac12e^{it}$ is too short to
change the winding of $e^{2it}$.) (c) If $h\colon\R\times S^1\to\R^2\setminus\{0\}$ is a smooth homotopy, Theorem~\ref{thm:homotopy} gives
$h_1^*\dd\theta-h_0^*\dd\theta=\dd K\dd\theta$, exact on $S^1$, so the integrals over $S^1$ agree.
\end{solution}

\begin{solution}{ex:toruscohomology}
\idea Each of the two independent loops on the torus is detected by the angle form that winds along it, and the area form is detected by its total integral.
\writeup (a) $\dd\theta_1$ is locally the differential of the function $\theta_1$, so it is closed; its integral over the circle $c_1(t)=(t,0)$, $t\in[0,1]$, is $1\ne0$, so it is not exact.
(b) If $a\,\dd\theta_1+b\,\dd\theta_2=\dd f$, integrating over $c_1$ and $c_2(t)=(0,t)$ gives $a=0$ and $b=0$. So the classes are independent and $\dim H^1(T^2)\ge2$. (c) Every $2$-form on
a surface is closed. If $\dd\theta_1\wedge\dd\theta_2=\dd\eta$, Stokes on the compact surface $T^2$ without boundary would give $1=\int_{T^2}\dd\theta_1\wedge\dd\theta_2=\int_{\emptyset}\eta=0$.
\end{solution}

\begin{solution}{ex:mobius}
\idea An orientation is a consistent choice of sign for top forms; the antipodal map reverses the sign on $S^2$, so no sign survives the passage to the quotient.
\writeup (a) $A^*x_i=-x_i$ and $A^*\dd x_i=-\dd x_i$, so each term of $\sigma$, of total degree $3$ in $x$ and $\dd x$, is multiplied by $(-1)^3$: $A^*\sigma=-\sigma$ on $\R^3$, hence on
$S^2$. (b) Let $\pi\colon S^2\to\R P^2$ be the quotient map, a local diffeomorphism (Exercise~\ref{ex:rp2}). If $\mu$ vanished nowhere on $\R P^2$, then $\pi^*\mu$ would vanish nowhere on $S^2$, so
$\pi^*\mu=f\sigma$ with $f$ smooth and never $0$. Since $\pi\circ A=\pi$, $A^*\pi^*\mu=\pi^*\mu$, that is, $(f\circ A)(-\sigma)=f\sigma$, so $f\circ A=-f$. But $S^2$ is connected and $f$ never vanishes,
so $f$ has constant sign, while $f(-x)=-f(x)$: contradiction. (c) On $S^n\subset\R^{n+1}$, $A^*\sigma=(-1)^{n+1}\sigma$; for $n$ odd this is $\sigma$, so $\sigma$ descends to a
nowhere-vanishing form on $\R P^n$. (For example $\R P^1$ is a circle, and $\R P^3=\mathrm{SO}(3)$, Chapter~7.)
\end{solution}

\begin{solution}{ex:spherevolume}
\idea Integrate the Gaussian in polar coordinates: the left side is known, the radial integral is a Gamma function, and what remains is the area of the sphere.
\writeup (a) Example~\ref{ex:sigma}. (b) In polar coordinates $x=r\xi$ ($r>0$, $\xi\in S^{n-1}$), Lebesgue measure is $r^{n-1}\dd r\,\dd S(\xi)$: the map
$(r,\xi)\mapsto r\xi$ pulls $\dd x^1\wedge\dots\wedge\dd x^n$ back to $r^{n-1}\dd r\wedge\sigma$, as one checks from $\dd(r\xi)=\xi\dd r+r\,\dd\xi$ and $\sigma_\xi(\dd\xi,\dots)$ being the area form.
So $\pi^{n/2}=s_{n-1}\int_0^\infty e^{-r^2}r^{n-1}\dd r=s_{n-1}\cdot\frac12\Gamma(\frac n2)$, by the substitution $u=r^2$. (c) $\omega_n=\frac{s_{n-1}}n=\frac{2\pi^{n/2}}{n\Gamma(n/2)}=\frac{\pi^{n/2}}{\Gamma(\frac n2+1)}$.
Since $\Gamma(\frac n2+1)$ grows faster than $C^n$ for every $C$ (Stirling), $\omega_n\to0$. The values $\omega_1,\omega_2,\dots$ are $2,\pi,\frac43\pi,\frac{\pi^2}2,\frac{8\pi^2}{15},\dots$, largest for
$n=5$ ($\approx5.26$). A ball of radius $1$ inside the cube $[-1,1]^n$ fills a fraction $\frac{\omega_n}{2^n}$, which is about $0.0025$ for $n=10$: in high dimension, almost
all of a cube's volume is in its corners, a fact that every algorithm sampling high-dimensional spaces has to live with.
\end{solution}

\section*{Sources and further reading}
\addcontentsline{toc}{section}{Sources and further reading}

V.~J.~Katz, ``The history of Stokes' theorem'', \emph{Math.\ Mag.} 52 (1979), 146--156, tells the story of Thomson's letter of 2 July 1850 and of the Smith's
Prize examination of February 1854; Thomson's letter is printed in \emph{The Correspondence between Sir George Gabriel Stokes and Sir William Thomson},
ed.\ D.~B.~Wilson (Cambridge, 1990). The examination paper is reproduced in Stokes' \emph{Mathematical and Physical Papers}, vol.~5 (Cambridge, 1905),
p.~320. For Maxwell's part, see C.~W.~F.~Everitt, \emph{James Clerk Maxwell: Physicist and Natural Philosopher} (Scribner, 1975).

\'E.~Cartan, ``Sur certaines expressions diff\'erentielles et le probl\`eme de Pfaff'', \emph{Ann.\ Sci.\ \'Ec.\ Norm.\ Sup.} (3) 16 (1899), 239--332;
G.~de~Rham, ``Sur l'analysis situs des vari\'et\'es \`a $n$ dimensions'', \emph{J.\ Math.\ Pures Appl.} (9) 10 (1931), 115--200. H.~Cartan,
\emph{Formes diff\'erentielles} (Hermann, 1967), is still the most elegant introduction to the calculus of forms.

Our proofs follow M.~Spivak, \emph{Calculus on Manifolds} (Benjamin, 1965), Chapters~4--5, whose preface warns that Stokes' theorem is ``trivial'' only
because the hard work is in the definitions, and J.~M.~Lee, \emph{Introduction to Smooth Manifolds}, 2nd ed., Chapters~14--17. R.~Bott and L.~W.~Tu,
\emph{Differential Forms in Algebraic Topology} (Springer, 1982), takes de Rham cohomology much further. The no-retraction proof of Brouwer's theorem by
integration goes back to the same circle of ideas as Milnor's article cited in Chapter~2; the ray construction of $r$ is standard.
